\documentclass[11pt, a4paper]{article}
\usepackage[utf8]{inputenc}
\usepackage[english]{babel}
\usepackage{graphicx}
\usepackage{float}
\usepackage{booktabs}
\usepackage{hyperref}
\usepackage{listings}
\usepackage{xcolor}
\usepackage{amsmath}
\usepackage{tikz}
\usepackage[left=3cm, right=3cm, top=2cm, bottom=3cm]{geometry}

\lstset{
    basicstyle=\ttfamily\footnotesize,
    breaklines=true,
    frame=single,
    backgroundcolor=\color{gray!10},
    rulecolor=\color{gray!30},
    keywordstyle=\color{blue},
    commentstyle=\color{green!50!black},
    stringstyle=\color{red},
    showstringspaces=false,
    tabsize=4
}

\title{
    \textbf{Software Quality 2025/26} \\[0.3cm]
    \large Project 2
}
\author{
    Eduardo Ervideira (72420) \\
    Noah Andriamampianina (76197)
}
\date{Maio 2026}

\begin{document}

\maketitle

\begin{center}
\textbf{GitHub repository:} \url{https://github.com/eduardoervideira/Bukkit} \\[0.2cm]
\end{center}

% \vspace{1cm}
% \tableofcontents
% \newpage

\section{Introduction}
This report documents the application of black-box and white-box testing techniques to Bukkit, which is an open source Java API and plugin framework for Minecraft multiplayer servers. 

Two functions were selected for testing:
\begin{itemize}
    \item \texttt{Color.fromRGB(int red, int green, int blue)}: Validates and creates a \texttt{Color} object from RGB values (0-255).
    \item \texttt{Note(int octave, Tone tone, boolean sharped)}: Validates and creates a musical note based on octave, tone, and sharpness.
\end{itemize}

The report is structured as follows:
\begin{itemize}
    \item Section \ref{sec:color} covers \texttt{Color.fromRGB} with Equivalence Class Testing, Boundary Value Analysis, and Control Flow-Based Testing.
    \item Section \ref{sec:note} covers \texttt{Note} with Cause-Effect Graphing, Category Partition Testing, and White-box Testing (Control Flow and Data Flow).
\end{itemize}

\section{Function: \texttt{Color.fromRGB}}
\label{sec:color}

\subsection{Source Code}
The \texttt{Color.fromRGB} function and its private constructor are shown below. The full implementation is available at: \href{[GITHUB_LINK_TO_COLOR_JAVA]}{\texttt{Color.java}}.

\begin{lstlisting}[language=Java, caption={Color.fromRGB Implementation}]
private Color(int red, int green, int blue) {
    Validate.isTrue(red >= 0 && red <= BIT_MASK, "Red is not between 0-255: ", red);
    Validate.isTrue(green >= 0 && green <= BIT_MASK, "Green is not between 0-255: ", green);
    Validate.isTrue(blue >= 0 && blue <= BIT_MASK, "Blue is not between 0-255: ", blue);

    this.red = (byte) red;
    this.green = (byte) green;
    this.blue = (byte) blue;
}

public static Color fromRGB(int red, int green, int blue) throws IllegalArgumentException {
    return new Color(red, green, blue);
}
\end{lstlisting}

\subsection{Equivalence Class Testing}
For rationale, the function validates three parameters (\texttt{red},
\texttt{green}, \texttt{blue}), each with a valid domain of [0, 255].
Equivalence Class Testing partitions the input into valid and invalid classes
for each parameter.

\subsubsection{Equivalence Classes}
\begin{table}[H]
\centering
\caption{Equivalence Classes for \texttt{Color.fromRGB}}
\label{tab:color_equivalence_classes}
\begin{tabular}{lccc}
\toprule
\textbf{Parameter} & \textbf{Invalid (Lower)} & \textbf{Valid} & \textbf{Invalid
(Upper)} \\
\midrule
\texttt{red}      & red $<$ 0   & 0 $\leq$ red $\leq$ 255   & red $>$ 255 \\
\texttt{green}    & green $<$ 0 & 0 $\leq$ green $\leq$ 255 & green $>$ 255 \\
\texttt{blue}     & blue $<$ 0  & 0 $\leq$ blue $\leq$ 255  & blue $>$ 255 \\
\bottomrule
\end{tabular}
\end{table}

Each parameter has 3 equivalence classes, therefore $|\texttt{red}| =
|\texttt{green}| = |\texttt{blue}| = 3$.

\subsubsection{Test Cases}
\textbf{Weak Equivalence:} The number of test cases equals the maximum number of
classes across all parameters: $\max(|\texttt{red}|, |\texttt{green}|,
|\texttt{blue}|) = 3$ tests.\\
\textbf{Strong Equivalence:} The cartesian product of partition subsets,
therefore $|\texttt{red}| * |\texttt{green}| * |\texttt{blue}| = 27$ tests.

\begin{table}[H]
\centering
\caption{Weak Equivalence Test Cases}
\label{tab:color_weak_equivalence}
\begin{tabular}{cccccc}
\toprule
\textbf{Test ID} & \textbf{Red} & \textbf{Green} & \textbf{Blue} &
\textbf{Classes Covered} & \textbf{Expected Outcome} \\
\midrule
WE1 & -1  & -1  & -1  & Invalid Lower (R, G, B) &
\texttt{IllegalArgumentException} \\
WE2 & 128 & 128 & 128 & Valid (R, G, B)         & \texttt{Color} object created
\\
WE3 & 256 & 256 & 256 & Invalid Upper (R, G, B) &
\texttt{IllegalArgumentException} \\
\bottomrule
\end{tabular}
\end{table}

\begin{table}[H]
\centering
\caption{Strong Equivalence Test Cases}
\label{tab:color_strong_equivalence}
\begin{tabular}{ccccc}
\toprule
\textbf{Test ID} & \textbf{Red} & \textbf{Green} & \textbf{Blue} &
\textbf{Expected Outcome} \\
\midrule
SE1  & -1  & -1  & -1  & \texttt{IllegalArgumentException} \\
SE2  & -1  & -1  & 128 & \texttt{IllegalArgumentException} \\
SE3  & -1  & -1  & 256 & \texttt{IllegalArgumentException} \\
SE4  & -1  & 128 & -1  & \texttt{IllegalArgumentException} \\
SE5  & -1  & 128 & 128 & \texttt{IllegalArgumentException} \\
SE6  & -1  & 128 & 256 & \texttt{IllegalArgumentException} \\
SE7  & -1  & 256 & -1  & \texttt{IllegalArgumentException} \\
SE8  & -1  & 256 & 128 & \texttt{IllegalArgumentException} \\
SE9  & -1  & 256 & 256 & \texttt{IllegalArgumentException} \\
SE10 & 128 & -1  & -1  & \texttt{IllegalArgumentException} \\
SE11 & 128 & -1  & 128 & \texttt{IllegalArgumentException} \\
SE12 & 128 & -1  & 256 & \texttt{IllegalArgumentException} \\
SE13 & 128 & 128 & -1  & \texttt{IllegalArgumentException} \\
SE14 & 128 & 128 & 128 & \texttt{Color} object created \\
SE15 & 128 & 128 & 256 & \texttt{IllegalArgumentException} \\
SE16 & 128 & 256 & -1  & \texttt{IllegalArgumentException} \\
SE17 & 128 & 256 & 128 & \texttt{IllegalArgumentException} \\
SE18 & 128 & 256 & 256 & \texttt{IllegalArgumentException} \\
SE19 & 256 & -1  & -1  & \texttt{IllegalArgumentException} \\
SE20 & 256 & -1  & 128 & \texttt{IllegalArgumentException} \\
SE21 & 256 & -1  & 256 & \texttt{IllegalArgumentException} \\
SE22 & 256 & 128 & -1  & \texttt{IllegalArgumentException} \\
SE23 & 256 & 128 & 128 & \texttt{IllegalArgumentException} \\
SE24 & 256 & 128 & 256 & \texttt{IllegalArgumentException} \\
SE25 & 256 & 256 & -1  & \texttt{IllegalArgumentException} \\
SE26 & 256 & 256 & 128 & \texttt{IllegalArgumentException} \\
SE27 & 256 & 256 & 256 & \texttt{IllegalArgumentException} \\
\bottomrule
\end{tabular}
\end{table}

\subsection{Boundary Value Analysis (BVA)}
Boundary Value Analysis focuses on testing values at and around the boundaries
of each input domain, where defects are statistically more likely to occur. The
boundary values considered are:
\begin{itemize}
    \item \textbf{Normal variants:} $\{0, 1, 128, 254, 255\}$ representing
    \textit{min}, \textit{min+}, \textit{nom}, \textit{max-}, and \textit{max}.
    \item \textbf{Robust variants:} additionally include $\{-1, 256\}$
    representing \textit{min-} and \textit{max+} (out-of-range values).
\end{itemize}

\subsubsection{Number of Test Cases}
\begin{table}[H]
\centering
\caption{Number of Test Cases for BVA Variants ($n = 3$)}
\label{tab:bva_summary}
\begin{tabular}{lcc}
\toprule
\textbf{Variant} & \textbf{Formula} & \textbf{Total} \\
\midrule
Simple BVA              & $4n + 1$ & 13  \\
Robust BVA              & $6n + 1$ & 19  \\
Worst Case BVA          & $5^n$    & 125 \\
Robust Worst Case BVA   & $7^n$    & 343 \\
\bottomrule
\end{tabular}
\end{table}

\paragraph{Rationale}
\begin{itemize}
    \item \textbf{Simple BVA} ($4n+1 = 13$): assumes the single-fault
    hypothesis. For each variable, the 5 boundary values are tested while
    keeping the remaining variables at their nominal value. The nominal-only
    case is counted once instead of $n$ times.
    \item \textbf{Robust BVA} ($6n+1 = 19$): extends Simple BVA by adding the
    two out-of-range values per variable, allowing explicit verification of the
    exception-throwing behavior.
    \item \textbf{Worst Case BVA} ($5^n = 125$): drops the single-fault
    hypothesis and considers all combinations of boundary values across the
    three parameters (Cartesian product).
    \item \textbf{Robust Worst Case BVA} ($7^n = 343$): the most exhaustive
    variant, combining the multi-fault assumption with the robust value set.
\end{itemize}

\subsubsection{Simple BVA Test Battery}
\begin{table}[H]
\centering
\caption{Simple BVA Test Cases for \texttt{Color.fromRGB}}
\label{tab:color_simple_bva}
\begin{tabular}{ccccc}
\toprule
\textbf{Test ID} & \textbf{Red} & \textbf{Green} & \textbf{Blue} &
\textbf{Expected Outcome} \\
\midrule
SBVA1  & 0   & 128 & 128 & \texttt{Color} object created \\
SBVA2  & 1   & 128 & 128 & \texttt{Color} object created \\
SBVA3  & 254 & 128 & 128 & \texttt{Color} object created \\
SBVA4  & 255 & 128 & 128 & \texttt{Color} object created \\
SBVA5  & 128 & 0   & 128 & \texttt{Color} object created \\
SBVA6  & 128 & 1   & 128 & \texttt{Color} object created \\
SBVA7  & 128 & 254 & 128 & \texttt{Color} object created \\
SBVA8  & 128 & 255 & 128 & \texttt{Color} object created \\
SBVA9  & 128 & 128 & 0   & \texttt{Color} object created \\
SBVA10 & 128 & 128 & 1   & \texttt{Color} object created \\
SBVA11 & 128 & 128 & 254 & \texttt{Color} object created \\
SBVA12 & 128 & 128 & 255 & \texttt{Color} object created \\
SBVA13 & 128 & 128 & 128 & \texttt{Color} object created \\
\bottomrule
\end{tabular}
\end{table}

\subsubsection{Robust BVA Test Battery}
\begin{table}[H]
\centering
\caption{Robust BVA Test Cases for \texttt{Color.fromRGB}}
\label{tab:color_robust_bva}
\begin{tabular}{ccccc}
\toprule
\textbf{Test ID} & \textbf{Red} & \textbf{Green} & \textbf{Blue} &
\textbf{Expected Outcome} \\
\midrule
RBVA1  & -1  & 128 & 128 & \texttt{IllegalArgumentException} \\
RBVA2  & 0   & 128 & 128 & \texttt{Color} object created \\
RBVA3  & 1   & 128 & 128 & \texttt{Color} object created \\
RBVA4  & 254 & 128 & 128 & \texttt{Color} object created \\
RBVA5  & 255 & 128 & 128 & \texttt{Color} object created \\
RBVA6  & 256 & 128 & 128 & \texttt{IllegalArgumentException} \\
RBVA7  & 128 & -1  & 128 & \texttt{IllegalArgumentException} \\
RBVA8  & 128 & 0   & 128 & \texttt{Color} object created \\
RBVA9  & 128 & 1   & 128 & \texttt{Color} object created \\
RBVA10 & 128 & 254 & 128 & \texttt{Color} object created \\
RBVA11 & 128 & 255 & 128 & \texttt{Color} object created \\
RBVA12 & 128 & 256 & 128 & \texttt{IllegalArgumentException} \\
RBVA13 & 128 & 128 & -1  & \texttt{IllegalArgumentException} \\
RBVA14 & 128 & 128 & 0   & \texttt{Color} object created \\
RBVA15 & 128 & 128 & 1   & \texttt{Color} object created \\
RBVA16 & 128 & 128 & 254 & \texttt{Color} object created \\
RBVA17 & 128 & 128 & 255 & \texttt{Color} object created \\
RBVA18 & 128 & 128 & 256 & \texttt{IllegalArgumentException} \\
RBVA19 & 128 & 128 & 128 & \texttt{Color} object created \\
\bottomrule
\end{tabular}
\end{table}

\subsubsection{Worst Case BVA Test Battery}
The Worst Case BVA test battery consists of all $5^3 = 125$ combinations
obtained from the Cartesian product of the boundary value set $\{0, 1, 128, 254,
255\}$ for the three parameters. Since all values lie within the valid domain
$[0, 255]$, \emph{all 125 test cases are expected to successfully create a
\texttt{Color} object}. This evidences a known limitation of this technique:
without including invalid values, Worst Case BVA does not exercise the
exception-throwing behavior of the function.

\begin{table}[H]
\centering
\caption{Worst Case BVA Test Cases for \texttt{Color.fromRGB} (representative
subset)}
\label{tab:color_worst_case_bva}
\begin{tabular}{ccccc}
\toprule
\textbf{Test ID} & \textbf{Red} & \textbf{Green} & \textbf{Blue} &
\textbf{Expected Outcome} \\
\midrule
WCBVA1   & 0   & 0   & 0   & \texttt{Color} object created \\
WCBVA2   & 0   & 0   & 1   & \texttt{Color} object created \\
WCBVA3   & 0   & 0   & 128 & \texttt{Color} object created \\
WCBVA4   & 0   & 0   & 254 & \texttt{Color} object created \\
WCBVA5   & 0   & 0   & 255 & \texttt{Color} object created \\
WCBVA6   & 0   & 1   & 0   & \texttt{Color} object created \\
\vdots   & \vdots & \vdots & \vdots & \vdots \\
WCBVA63  & 128 & 128 & 128 & \texttt{Color} object created \\
\vdots   & \vdots & \vdots & \vdots & \vdots \\
WCBVA125 & 255 & 255 & 255 & \texttt{Color} object created \\
\bottomrule
\end{tabular}
\end{table}

\subsubsection{Robust Worst Case BVA Test Battery}
The Robust Worst Case BVA test battery consists of all $7^3 = 343$ combinations
obtained from the Cartesian product of the robust boundary value set $\{-1, 0,
1, 128, 254, 255, 256\}$ for the three parameters. A test case produces a valid
\texttt{Color} object only if all three parameters are within $[0, 255]$;
otherwise, an \texttt{IllegalArgumentException} is thrown. This yields $5^3 =
125$ successful cases and $343 - 125 = 218$ exception cases.

\begin{table}[H]
\centering
\caption{Robust Worst Case BVA Test Cases for \texttt{Color.fromRGB}
(representative subset)}
\label{tab:color_robust_worst_case_bva}
\begin{tabular}{ccccc}
\toprule
\textbf{Test ID} & \textbf{Red} & \textbf{Green} & \textbf{Blue} &
\textbf{Expected Outcome} \\
\midrule
RWCBVA1   & -1  & -1  & -1  & \texttt{IllegalArgumentException} \\
RWCBVA2   & -1  & -1  & 0   & \texttt{IllegalArgumentException} \\
RWCBVA3   & -1  & -1  & 1   & \texttt{IllegalArgumentException} \\
\vdots    & \vdots & \vdots & \vdots & \vdots \\
RWCBVA172 & 128 & 128 & 128 & \texttt{Color} object created \\
\vdots    & \vdots & \vdots & \vdots & \vdots \\
RWCBVA343 & 256 & 256 & 256 & \texttt{IllegalArgumentException} \\
\bottomrule
\end{tabular}
\end{table}

\subsection{Used Prompt}
\begin{lstlisting}
I have a Java method Color.fromRGB(int red, int green, int blue) that validates
each parameter to be in [0, 255] and throws IllegalArgumentException otherwise.
Generate JUnit 5 tests from these specifications:

Simple BVA (13 tests, SBVA1-SBVA13): vary one parameter at a time across {0, 1,
128, 254, 255}, keep others at 128. 
Robust BVA (19 tests): Simple BVA + out-of-range values -1 and 256 per
parameter. 
Worst Case BVA (125 tests): Cartesian product of {0, 1, 128, 254, 255}. 
Robust Worst Case BVA (343 tests): Cartesian product of {-1, 0, 1, 128, 254,
255, 256}.

Use @ParameterizedTest with @CsvSource for the smaller batteries and
@MethodSource with programmatic generation for the larger ones. Each test
asserts either successful Color creation (all parameters in [0, 255]) or
IllegalArgumentException.
\end{lstlisting}

\subsection{Control Flow-Based Testing}
For rationale, the function has three validation checks
(\texttt{Validate.isTrue}) that throw exceptions if any condition fails. Control
Flow-Based Testing ensures all statements, decisions, and paths are covered.

\textbf{TODO}

\section{Function: \texttt{Note}}
\label{sec:note}
Placeholder for future content.

\section{Conclusion}
Placeholder for future content.

\end{document}